\section{Contribution to the Dissertation}
\label{sec:otel-weaver-exp-005-collector-to-intake-contribution}

\subsection{Contribution to Prediction vs.~Coordination}
\label{sec:otel-weaver-exp-005-collector-to-intake-contribution-prediction}

The central thesis claim is that process intelligence requires coordination evidence, not
prediction output. EXP-005 contributes the first operational demonstration that the boundary
between prediction-layer telemetry and coordination-layer process evidence can be enforced
at the infrastructure level without application-side changes. The \texttt{filter/pi} OTTL
processor is the operational instantiation of this boundary: it enforces the presence of
\texttt{process.pi.instance\_id}, which is the attribute that binds a span to a coordination
event (an instance of a named process) rather than treating it as a prediction-layer
observation (a metric from an arbitrary service). The distinction is not theoretical; it is
enforced by a silently-dropping filter that is architecturally prior to any process mining
stage.

This contribution is PARTIAL in its evidentiary weight: the architecture is correct and the
specification is complete, but the claim that the filter \textit{actually drops} Span~B
when connected to a live OTLP stream has not been confirmed by observed output. The
contribution is therefore admitted as a design-level contribution, not an experimental-result
contribution, until the live run is executed.

\subsection{Contribution to Process Evidence}
\label{sec:otel-weaver-exp-005-collector-to-intake-contribution-process-evidence}

EXP-005 contributes the OCEL~2.0 JSON output schema as the process evidence format for the
OTel Weaver pipeline series. The \texttt{transform/pi} processor's injection of
\texttt{pi.intake.ingested\_at} (collector-clock ISO~8601 timestamp) is a specific
contribution to the process evidence quality argument: by anchoring the event timestamp to
the collector clock rather than the emitter clock, the experiment eliminates clock-skew
between emitting services as a source of process ordering uncertainty. This is a direct
requirement of the Van der Aalst Constitution (temporal conformance criterion) from the
process-mining-chicago-tdd rules: stages must occur in lawful order with no impossible
overlaps, which requires a single authoritative clock for event timestamps in the log.

The three-class Span~A/B/C taxonomy is a process evidence quality taxonomy: Span~A
represents evidence that can be mined, Span~B represents non-evidence that must be excluded,
and Span~C represents evidence that can be mined but will fail conformance court validation.
This taxonomy is the first classification of live OTLP spans by their process evidence quality,
and it provides the dissertation with a concrete, falsifiable test surface for the claim that
process evidence can be separated from non-evidence at the collector boundary.

\subsection{Contribution to Manufacturing Architecture}
\label{sec:otel-weaver-exp-005-collector-to-intake-contribution-manufacturing}

EXP-005 establishes the collector pipeline as an operator in the CodeManufactory manufacturing
architecture. Specifically, the three-stage pipeline (OTLP receiver $\to$ filter+transform
processors $\to$ file exporter) is a manufacturing pipeline in the canonical sense: it
receives feedstock (raw OTLP spans), applies lawful operators (filter by attribute presence,
transform by OTTL rules), and emits artifacts (OCEL~2.0 JSON records). The
\texttt{transform/pi} processor is a manufacturing operator in the exact sense of the
CodeManufactory terminology: it transforms feedstock into a canonical artifact form without
altering the upstream service's output.

The experiment also establishes the \texttt{BridgeRx} schema-version translation bridge as
a manufacturing primitive: it protects the conformance court from Weaver schema diffs by
translating span attributes to the canonical court-internal representation before validation.
This is an instance of the general principle that manufacturing pipelines should be
schema-version-isolated from both upstream feedstock and downstream courts, with explicit
bridge operators at each boundary. \textit{[AUTHOR THESIS: the BridgeRx bridge is referenced
in the collector-integration-model.yaml but its implementation is located in EXP-003's
validator.rs; its full implementation surface is outside the direct scope of EXP-005.]}

\subsection{Contribution to Capital-Flow Infrastructure}
\label{sec:otel-weaver-exp-005-collector-to-intake-contribution-capital-flow}

In the dissertation's capital-flow / settlement / DAO architecture layer, process intelligence
is the evidence layer that validates whether a capital-flow event (a transaction, a settlement,
a DAO governance vote) was executed by a lawful process or by an unlawful shortcut. EXP-005
contributes the live telemetry ingestion boundary for this evidence layer. A capital-flow
system instrumented with OpenTelemetry would emit spans for every significant transaction
event. EXP-005 demonstrates that those spans can be routed through the collector pipeline
and converted into OCEL~2.0 process evidence without requiring the capital-flow system to
implement process intelligence awareness. The collector boundary is the sovereignty boundary:
it is where the capital-flow system's observability output becomes the process intelligence
system's evidence input. The \texttt{process.pi.instance\_id} attribute is the capital-flow
process instance binding --- the identifier that links a span to a specific settlement
instance, DAO proposal, or transaction workflow. \textit{[AUTHOR THESIS: the capital-flow
binding is an architectural projection; EXP-005 does not implement or test a capital-flow
use case directly.]}
